%%%%%%%%%%%%%%%%%%%%%%%%%%%%%%%%%%%%%%%%%%%%%%%%%%%%%%%%%%
%%
%%  LOD-L1.tex
%%  Lecture 1: At the Beginning...
%%
%%%%%%%%%%%%%%%%%%%%%%%%%%%%%%%%%%%%%%%%%%%%%%%%%%%%%%%%%%

\cleardoublepage

\chapter*{\S L1. At the Beginning\dots}
\setcounter{footnote}{0}

%%%%%%%%%%%%%%%%%%%%%%%%%%%%%%%%%%%%%%%%%%%%%%%%%%%%%%%%%%
%%
%% MARK: Abstract
%%
%%%%%%%%%%%%%%%%%%%%%%%%%%%%%%%%%%%%%%%%%%%%%%%%%%%%%%%%%%

\begin{abstract}
  This lecture presents the prerequisites for the study of diffeology.
  As we shall see,
  a good understanding of infinitely differentiable maps on Euclidean domains,
  and a few fundamental theorems of calculus,
  will be a good start.
\end{abstract}

%%%%%%%%%%%%%%%%%%%%%%%%%%%%%%%%%%%%%%%%%%%%%%%%%%%%%%%%%%
%%
%% MARK: Introduction
%%
%%%%%%%%%%%%%%%%%%%%%%%%%%%%%%%%%%%%%%%%%%%%%%%%%%%%%%%%%%

Let us introduce some vocabulary.
First of all,
we call \emph{Euclidean space} any real vector space $\RR^n$ for some $n$.

The Euclidean structure of $\RR^n$,
defined by the inner product
$$%
  x \cdot y = \sum_{i=1}^n x^i y^i
  \quad \text{with} \quad
  \Norm{x} = \sqrt{\sum_{i=1}^n(x^i)^2},
  $$%
is used to define its (standard) topology:
an open subset $\cO$ in $\RR^n$ is any \emph{union of open balls},
like:
$$%
  \cO = \bigcup_{i \in \cI} \cB(r_i,\epsilon_i),
  $$%
where $\cI$ is any set of indices,
$r_i$ is any point in $\RR^n$,
and $\epsilon_i$ is any (strictly) positive number.

Therefore,
a subset $U \in \RR^n$ is open for the standard topology if (and only if):
for each $r \in U$ there exists an open ball $\cB(r,\epsilon)$ included in $U$,
$\cB(r,\epsilon) \subset U$.

Next,
we call \emph{Euclidean domain} any open subset of a Euclidean space.
We generally denote them by a capital letter $U, V, W$, etc.,
and also by cursive letters $\cO$, etc.

We also denote by $\Domains(\RR^n)$ the set of all Euclidean domains in $\RR^n$.
We call them simply \emph{$n$-domains}.

Note that this is just the topology of $\RR^n$,
sometimes denoted by $\Top(\RR^n)$.

A domain in $\RR^n$ is said to be of dimension $n$,
or be an $n$-domain.

Now,
about \emph{continuity}.
A map $F \colon U \to V$,
where $U$ and $V$ are Euclidean domains,
is said to be \emph{continuous} if (and only if):
the pullback $f^{-1}(\cO)$ of any open subset $\cO \subset V$ is an open subset in $U$.

That is equivalent to saying that for any open ball $\cB \subset V$,
there exists a family of open balls $\cB_i \subset U$ such that $f^{-1}(\cB) = \cup_i \cB_i$,
for some family of indices $\cI$.

%%%%%%%%%%%%%%%%%%%%%%%%%%%%%%%%%%%%%%%%%%%%%%%%%%%%%%%%%%
%%
%% MARK: Differentiable and smooth paths
%%
%%%%%%%%%%%%%%%%%%%%%%%%%%%%%%%%%%%%%%%%%%%%%%%%%%%%%%%%%%

\section{Differentiable and smooth paths}

Consider a \emph{path}
$$%
  \gamma \colon \OpenInterval{a,b} \to \RR^n.
  $$%

Assume that $\gamma$ is continuous.
Then,
consider the map
$$%
  \Delta\gamma \colon (t,t') \mapsto \frac{\gamma(t') - \gamma(t)}{t'-t}, \\
  $$%
defined on
$$%
  \OpenInterval{a,b}^2 \ - \ \{(t,t)\}_{a<t<b},
  $$%
with values in $\RR^n$.
The closed subset $\{(t,t)\}_{a<t<b}$ is called the \emph{diagonal}.
The function $\Delta \gamma$ is continuous,
but:

\begin{Article}{Definition.}{\itshape}
  If $\Delta\gamma$ extends continuously on the diagonal,
  we say that $\gamma$ is \smuline{continuously differentiable} or of \underline{class $\Ck{1}$}.
\end{Article}

We denote by
$$%
  \dot\gamma(t)
  \quad \text{or} \quad
  \frac{d\gamma(t)}{dt}
  \quad \text{the value} \quad \Delta\gamma(t,t) = \lim_{t' \to t} \Delta\gamma(t',t).
  $$%
The function $\dot\gamma \colon \OpenInterval{a,b} \to \RR^n$ is then called the \emph{derivative} of $\gamma$.
The value $\dot\gamma(t)$ is the derivative of $\gamma$ at the point $t$,
or {\itshape at time} $t$.

Now,
if $\gamma$ is of class $\Ck{1}$,
then $\dot\gamma$,
which is defined on the same interval,
is continuous.
One can ask if $\dot\gamma$ is also $\Ck{1}$?
If this is the case,
we generally denote by $\ddot\gamma$ the derivative of $\dot\gamma$.
We say that $\gamma$ is of class $\Ck{2}$.

That leads to the definition:

\begin{Article}{Definition.}{\itshape}
  We say that $\gamma$ is of \underline{class $\Ck{k}$},
  $k>1$,
  if (and only if):
  $$%
    \gamma \text{ is of class } \Ck{k-1} \text{ and } \frac{d^{k-1}\gamma(t)}{dt^{k-1}} \text{ is of class } \Ck{1}.
    $$%
  And we say that $\gamma$ is \underline{smooth},
  or of \underline{class $\Cinfty$},
  if $\gamma$ is of class $\Ck{k}$ for all $k \in \NN$.
\end{Article}
%
Note that continuous paths are said to be of \underline{class $\Czero$}.

%%%%%%%%%%%%%%%%%%%%%%%%%%%%%%%%%%%%%%%%%%%%%%%%%%%%%%%%%%
%%
%% MARK: Smooth Maps, the Holistic Approach
%%
%%%%%%%%%%%%%%%%%%%%%%%%%%%%%%%%%%%%%%%%%%%%%%%%%%%%%%%%%%

\section{Smooth maps, the holistic approach}

In diffeology we are interested principally in {\itshape smooth maps},
or {\itshape infinitely differentiable maps}.
Precisely,
diffeology consists in extending the concept of smooth maps from Euclidean domains to arbitrary sets.
That is why we have the choice in the definitions of smooth maps between Euclidean domains.
We will present them here,
not in their historic order but according to what we feel is more intuitive.
The first definition is actually a theorem.

\begin{Article}{Definition-Theorem.}{\itshape}{\upshape (Boman 1967)}\label{Boman-theorem}
  Let $F \colon U \to V$ be a continuous map between two Euclidean domains.
  The map $F$ is \smuline{infinitely diffe}- \smuline{rentiable}, %% PIZ Adjustement
  or \underline{smooth},
  if (and only if) the composite $F \circ \gamma$,
  where $\gamma$ is any smooth path in $U$,
  is a smooth path in $V$.
\end{Article}

As we have seen,
smooth paths have been defined previously,
independently.

%%%%%%%%%%%%%%%%%%%%%%%%%%%%%%%%%%%%%%%%%%%%%%%%%%%%%%%%%%
%%
%% MARK: Differentiable Maps, the Pedestrian Approach
%%
%%%%%%%%%%%%%%%%%%%%%%%%%%%%%%%%%%%%%%%%%%%%%%%%%%%%%%%%%%

\section{Differentiable maps, the pedestrian approach}

Now,
the classical definition of differentiable map begins with the definition of \emph{tangent maps}.

\begin{Article}{Definition.}{\itshape}
  Let $f$ and $g$ be two maps defined on a same Euclidean domain $U$ to some $\RR^m$.
  Let $r_0$ be some point in $U$.
  We say that $f$ and $g$ are \smuline{tangent} in $r_0$ if:
  $$%
    \lim_{r \to r_0} \frac{f(r) - g(r)}{\Norm{r-r_0}} = 0.
    $$%
  Of course,
  that implies in particular that $f(r_0) = g(r_0)$.
\end{Article}

\begin{Article}{Definition.}{\itshape}
  Let $f$ be a map defined on some domain $U \in \RR^n$ with values in $\RR^m$.
  We say that $f$ is \underline{differentiable},
  or \underline{derivable},%
  \footnote{The words {\itshape differentiable} and {\itshape derivable} are here completely equivalent.}
  at the point $r_0 \in U$ if it is tangent,
  at this point,
  to an \emph{affine map}
  $$%
    r \mapsto f(r_0) + M(r-r_0),
    $$%
  where $M \in \Lin(\RR^n,\RR^m)$ is a linear map from $\RR^n$ to $\RR^m$.
  That condition means precisely that,
  \begin{equation}\tag{$\clubsuit$}
    \lim_{r \to r_0} \frac{f(r)- f(r_0)- M(r-r_0)}{\Norm{r-r_0}} = 0.
  \end{equation}
  We say that $f$ is \emph{differentiable on $U$},
  or simply differentiable,
  without more precision,
  if $f$ is tangent to an affine map everywhere.
\end{Article}

\begin{Article}{Remark.}{}\label{continuously-differentiable}
  The condition ($\clubsuit$) implies that $f(r)$ converges to $f(r_0)$ when $r$ tends to $r_0$.
  That is,
  $f$ is continuous in $r_0$:
  to be tangent to an affine map implies to be continuous:
  \emph{a differentiable map is necessarily continuous}.
\end{Article}

\begin{proof}
  Since $\lim_{r \to r_0} \frac{f(r)- f(r_0)- M(r-r_0)}{\Norm{r-r_0}} = 0$,
  $\lim_{r \to r_0} (f(r)- f(r_0)- M(r-r_0)) = 0$.
  Thus,
  since $\lim_{r \to r_0}M(r-r_0)=0$,
  $\lim_{r \to r_0} (f(r)- f(r_0)) = 0$.
  That is $f(r) \xrightarrow[r \to r_0]{} f(r_0)$,
  and $f$ is continuous at $r_0$.
  That is the logic of the situation.
  More formally,
  %since $\lim_{r \to r_0} \frac{f(r)- f(r_0)- M(r-r_0)}{\Norm{r-r_0}} = 0$,
  since $\lim_{r \to r_0} (f(r)- f(r_0)- M(r-r_0))/\Norm{r-r_0} = 0$,
  for all $\epsilon > 0$,
  there exists $\eta >0$ such that $\Norm{r-r_0}<\eta$ implies $\Norm{f(r)- f(r_0) - M(r-r_0)} < \epsilon \Norm{r-r_0} < \epsilon \eta$.
  On the other hand,
  let $m$ be $\Norm{M} = \Sup_{\Norm{u}=1} \Norm{M(u)}$ (remember, the sphere $\Sph^{n-1}$ is compact).
  Thus,
  for all $r$,
  $\Norm{M(r-r_0)} \leq m \Norm{r-r_0}$.
  That is,
  for all $\epsilon >0$ there is $\eta'>0$ such that $\Norm{r-r_0}<\eta'$ implies $\Norm{M(r-r_0)} < \epsilon$;
  actually $\eta'$ can be chosen equal to $\epsilon/m$.
  Now,
  the inequality on a triangle $u + v = w$ says that $\Norm{w} \leq \Norm{u} + \Norm{v}$,
  or $\Norm{w} - \Norm{u} \leq \Norm{v}$.
  Applied to $u = M(r-r_0)$,
  $w = f(r)-f(r_0)$ and $v = w-u = f(r)-f(r_0) - M(r-r_0)$,
  that gives:
  $\Norm{f(r)-f(r_0)} - \Norm{M(r-r_0)} \leq \Norm{f(r)- f(r_0) - M(r-r_0)} \leq \epsilon$.
  Then,
  $\Norm{f(r)-f(r_0)} \leq \epsilon + \Norm{M(r-r_0)}$.
  Now,
  choosing $\eta'' < \inf(\eta, \eta')$,
  we get $\Norm{f(r)-f(r_0)} \leq \epsilon + \epsilon$.
  Changing $\epsilon$ to $\epsilon/2$,
  we get:
  for all $\epsilon > 0$,
  there exists $\eta'' > 0$ such that $\Norm{r-r_0} < \eta''$ implies $\Norm{f(r)-f(r_0)} \leq \epsilon$.
  Therefore,
  $f$ is continuous at $r_0$.
\end{proof}

Now let us describe more precisely the linear part of the affine tangent map.
Pick a vector $v \in \RR^n$,
$v \neq 0$,
and let $r = r_0 + tv$,
where $t >0$ is small enough for $r$ to be in the domain $U$.
The last condition writes:
$$%
  \lim_{t \to 0} {f(r_0 + tv)- f(r_0)- M(tv) \over t \Norm{v}} = 0.
  $$%
And then,
$$%
  M(v) = \lim_{t \to 0} {f(r_0 + tv)- f(r_0) \over t}.
  $$%
The linear map $M$ is then called the \emph{tangent linear map} and it is denoted by
$$%
  \Der(f)_{r_0}(v) \text{ or } \Der(f)(r_0)(v) \text{ or } \Der(f)(r_0)(v) = \lim_{t \to 0} {f(r_0 + tv)- f(r_0) \over t}.
  $$%
So,
$f$ is differentiable on $U$ if it admits a tangent linear map at every point in $U$,
$$%
  f' \text{ or } \Der(f) \text{ or } \Der(F) \colon U \to \Lin(\RR^n,\RR^m).
  $$%
The \emph{affine tangent map} at the point $r_0$ then writes
$$%
  r \mapsto f(r_0) + \Der(f)(r_0)(r-r_0).
  $$%
Note that there is another way to express the approximation of the function $f$ by its affine tangent map around $r_0$:
$$%
  f(r) = f(r_0) + \Der(f)(r_0)(r-r_0) + o(\Norm{r-r_0}),
  $$%
where Landau's \emph{Little-O} notation $o(x)$ means $\underset{x \to 0}{\lim} o(x)/x = 0$.

As we know,
the set of linear maps $\Lin(\RR^n,\RR^m)$ is a real vector space of dimension $n \times m$,
equivalent to $\RR^{n\times m}$.
As such:

\begin{Article}{Definition.}{\itshape}
  We say that $f \colon U \to \RR^m$ is of \underline{class $\Ck{1}$} if $f$ is differentiable on $U$,
  and if the map that associates its tangent linear map with each point in $U$ is continuous.
  That is,
  if the map $r \mapsto \Der(f)(r)$,
  defined on $U$ with values in $\Lin(\RR^n,\RR^m)$,
  is continuous.
\end{Article}

\begin{Article}{Remark.}{}
  An affine map from $\RR^n$ to $\RR^m$ writes $x \mapsto Ax + b$,
  where $A \in \Lin(\RR^n,\RR^m)$ and $b \in \RR^m$.
  The set of these affine maps is denoted by $\Aff(\RR^n,\RR^m)$;
  it is naturally a vector space of dimension $m + n \times m$,
  equivalent to the space of matrices
  $$%
    \Vector{A & b \\ 0 & 1}
    \quad \text{such that} \quad
    \Vector{A & b \\ 0 & 1} \Vector{x \\ 1} = \Vector{A x + b \\ 1},
    $$%
  and inherits the standard topology of $\RR^{m + n \times m}$.
  The affine tangent map writes, for all $r \in U$:
  $$%
    \Vector{\Der(f)_{r} & f(r) - \Der(f)_{r}(r) \\ 0 & 1}.
    $$%
  We can notice that,
  to be $\Ck{1}$,
  it is equivalent to request that the maps that associate the linear tangent map in $\Lin(\RR^n,\RR^m)$,
  or the affine tangent map in $\Aff(\RR^n,\RR^m)$,
  with every point $r$ in $U$,
  are continuous.
  However,
  on a conceptual level:
  to be differentiable implies to be continuous;
  then,
  having a \emph{first local affine approximation}
  ---~which is exactly what the affine tangent map is~---
  continuous everywhere is a natural request.
  And that is the meaning of being $\Ck{1}$.
\end{Article}

%%%%%%%%%%%%%%%%%%%%%%%%%%%%%%%%%%%%%%%%%%%%%%%%%%%%%%%%%%
%%
%% MARK: Higher Order Derivatives and Smooth Maps
%%
%%%%%%%%%%%%%%%%%%%%%%%%%%%%%%%%%%%%%%%%%%%%%%%%%%%%%%%%%%

\section{Higher order derivatives and smooth maps}

Now,
we can look for \emph{higher order derivatives}.
Let us begin by the second order.
Assume that $f$ is derivable,
then its derivative $\Der(F)$ is defined on $U$ with values in $\Lin(\RR^n,\RR^m)$.
The second derivative of $f$ is then defined by
$$%
  \Der^2(f)(r) = \Der\big(r \mapsto \Der(f)(r)\big)(r) \in \L\big(\RR^n, \Lin(\RR^n, \RR^m)\big).
  $$%
Note that if $f$ is derivable and if its derivative $\Der(f)$ is derivable,
that implies that $\Der(f)$ is continuous and $f$ is $\Ck{1}$.

So,
we get a recursive definition of higher derivatives of $f$ and then a recursive definition of class $\Ck{k}$:

\begin{Article}{Definition.}{\itshape}\label{Higher-derivatives}
  The $k$th-derivative of $f$,
  if it exists,
  is the derivative denoted and defined recursively by
  $$%
    \Der^k(f) = \Der\big(r \mapsto \Der^{k-1}(f)(r)\big).
    $$%
\end{Article}
%
Note that,
if $f$ is differentiable up to order $k$,
then $\Der^{\ell}(f)$ is continuous up to $\ell = k-1$,
thanks to Remark \ref{continuously-differentiable}.
Therefore,
if $f$ admits a $k$th-derivative,
then $f$ is $\Ck{k-1}$,
with the definition:

\begin{Article}{Definition.}{\itshape}
  A \emph{function $f$ is $\Ck{k}$},
  or of \emph{class $\Ck{k}$},
  if $f$ is derivable up to order $k$ and its $k$-derivative is continuous.
  
  We say that $f$ is \emph{infinitely differentiable},
  or \emph{infinitely derivable},
  or \emph{smooth},
  if $f$ is of class $\Ck{k}$ for all integer $k$.
\end{Article}

\begin{Article}{Remark.}{}
  Boman's theorem cited previously says exactly that Definition \ref{Higher-derivatives} and Definition-Theorem \ref{Boman-theorem} are coherent.
  Beware,
  it is a subtle and non-trivial theorem.
\end{Article}

%%%%%%%%%%%%%%%%%%%%%%%%%%%%%%%%%%%%%%%%%%%%%%%%%%%%%%%%%%
%%
%% MARK: The Tangent Linear Map
%%
%%%%%%%%%%%%%%%%%%%%%%%%%%%%%%%%%%%%%%%%%%%%%%%%%%%%%%%%%%

\section{The tangent linear map}

We consider a differentiable map $f \colon U \to \RR^m$,
with $U \subset \RR^n$ an open subset.

Let us denote by $x$ and $y$ the source and target variables involved in $f$,
that is,
$f \colon x \mapsto y$ with $x = (x_1,\dots,x_n) \in U$ and $y = (y_1,\dots,y_m) \in \RR^m$.
The tangent linear map can be written indifferently
$$%
  \Der(f)(x) \text{ or } \Der(x \mapsto y)(x),
  $$%
depending on what we want to focus on.
For example,
if you denote the square root function by $\SquareRoot$,
you can write its derivative as $\Der(\SquareRoot)(x)$,
or $\Der(x \mapsto \sqrt{x}\,)(x)$.

Now,
decompose $x$, $v$ and $y$ on the canonical basis,
$$%
  x = \sum_{i=1}^n x^i \bme_i,
  \quad  v = \sum_{j=1}^m v^i \bme_i
  \quad \text{and} \quad
  y = \sum_{j=1}^m y^j \bme_j
  \quad \text{with} \quad y^j = f^j(x).
  $$%
Now,
the tangent linear map writes
\begin{eqnarray*}\allowdisplaybreaks
  \Der(x \mapsto y)(x)(v) & = & \Der(x \mapsto y)(x)\left(\sum_{i=1}^n v^i \bme_i\right) \\
  & = & \sum_{i=1}^n v^i \Der(x \mapsto y)(x)(\bme_i) \\
  & = & \sum_{i=1}^n v^i \Der\left(x \mapsto \sum_{j=1}^m y^j \bme_j\right)(x)(\bme_i)\\
  & = & \sum_{i=1}^n \sum_{j=1}^m v^i \Der\left(x \mapsto y^j \right)(x)(\bme_i) \cdot \bme_j
\end{eqnarray*}
In other words,
the tangent linear map $\Der(x \mapsto y)(x)$ is represented by the matrix $(\Der_i^j)_{i = 1}^n{}_{j=1}^m$,
where
$$%
  \Der_i^j = \Der\left(x \mapsto y^j \right)(x)(\bme_i).
  $$%
Now,
what does $\Der\left(x \mapsto y^j \right)(x)(\bme_i)$ represent exactly?
$$%
  \Der\left(x \mapsto y^j \right)(x)(\bme_i) =
  \lim_{t\to 0} {f^j(x_1,\ldots,x_i + t,\ldots,x_n) - f^j(x_1,\ldots,x_i,\ldots,x_n) \over t},
  $$%
which is,
by definition,
the partial derivative
$$%
  \Der_i^j = {\partial y^j \over \partial x^i}
  \quad \text{also denoted by} \quad \partial_i y^j.
  $$%

\underline{Notations:} For convenience we will denote also
$$%
  \begin{array}{lll}
  \Der(x\mapsto y)(x) & \text{by}  &  {\partial y \over \partial x}, \text{ the derivative}; \\[0.1cm]
  \text{and}        &            & \\[0.0cm]
  \Der(x\mapsto y)(a) & \text{by} &  \left.{\partial y \over \partial x}\right|_{x=a}, \text{ the value of the derivative}.
  \end{array}
  $$%
Eventually,
the tangent linear map can be written as the matrix of partial derivatives
$$%
  \Der(x\mapsto y)(x) = \Der\left(\Vector{x^1 \\ \vdots \\ x^n} \mapsto \Vector{y^1 \\ \vdots \\ y^m}\right)(x)
  = \Vector{\partial_1 y^1 & \cdots & \partial_n y^1 \\
  \vdots & \ddots & \vdots \\
  \partial_1 y^m & \cdots & \partial_n y^m
  }
  $$%
such that
$$%
  \Der(x\mapsto y)(x)(v) = \Vector{\partial_1 y^1 & \cdots & \partial_n y^1 \\
  \vdots & \ddots & \vdots \\
  \partial_1 y^m & \cdots & \partial_n y^m
  }\Vector{v^1 \\ \vdots \\ v^n}
  = \Vector{\sum_{i=1}^n v^i \partial_i y^1 \\ \vdots \\ \sum_{i=1}^n v^i \partial_i y^m}.
  $$%
\begin{Article}{Proposition.}{\itshape}
  Let $f \colon  U \to V$ and $g \colon V \to W$ be two $\Ck{1}$ maps,
  with $U \subset \RR^n$,
  $V \subset \RR^m$ and $W \subset \RR^\ell$.
  The derivative of the composite,
  which is also $\Ck{1}$,
  satisfies the so-called \emph{chain-rule}:
  $$%
    \Der(g \circ f)(r) = \Der(g)(f(r)) \circ \Der(f)(r).
    $$%
  Note also that the derivation $\Der$ is linear.
  If $f_1$ and $f_2$ are defined from $U$ to $V$ and $\lambda, \mu \in \RR$,
  then:
  $$%
    \Der(\lambda f_1 + \mu f_2)) = \lambda \Der(f_1) + \mu \Der(f_2).
    $$%
\end{Article}
%
\begin{Article}{Note.}{}
  The chain-rule writes in terms of partial derivatives:
  $$%
    {\partial z \over \partial x} = {\partial z \over \partial y} {\partial y \over \partial x},
    $$%
  and the value at the point $x=a$,
  where $b = f(a)$:
  $$%
    \left.{\partial z \over \partial x}\right|_{x=a} =
    \left.{\partial z \over \partial y}\right|_{y=b} \left.{\partial y \over \partial x}\right|_{x=a}.
    $$%
  In terms of coordinates:
  $$%
    \partial_i z^j = \sum_{k=1}^m \partial_i y^k \partial_k z^j
    \quad \text{that is} \quad
    {\partial z^j \over \partial x^i} = \sum_{k=1}^m {\partial z^j \over \partial y^k} {\partial y^k \over \partial x^i},
    $$%
  where $i = 1 \dots n$,
  $k = 1 \dots m$ and $j = 1 \dots \ell$.
\end{Article}
%
\begin{Article}{Note.}{}
  In particular,
  for $\gamma \colon \cI \to U$,
  a smooth path in $U$,
  $$%
    \Der(f \circ \gamma)(t)(1) = \Der(f)(\gamma(t)) \circ \Der(\gamma)(t)(1),
    $$%
  but
  $$%
    \Der(\gamma)(t)(1) = \lim_{\epsilon \to 0} {\gamma(t+\epsilon \times 1) - \gamma(t) \over \epsilon} = {d\gamma(t) \over dt} = \dot\gamma(t).
    $$%
  Therefore:
  $$%
    \Der(f \circ \gamma)(t)(1) = \Der(f)(\gamma(t))\big(\dot\gamma(t)\big).
    $$%
\end{Article}
%
\begin{Article}{Notation.}{}
  It happens that a tangent vector $v \in \RR^n$ would be denoted by a variational notation $\delta x$,
  meaning that $v$ is regarded as a \emph{variation} of the point $x$.
  Write $\gamma \colon t \mapsto x$ and $f \colon x \mapsto y$,
  then $f \circ \gamma \colon t \mapsto y$,
  meaning $t \mapsto x \mapsto y$.
  Using the partial derivative notation introduced previously,
  the formula above writes:
  $$%
    \delta y = {\partial y \over \partial x}(\delta x),
    $$%
  with
  $$%
    x = \gamma(t),\ y = f(x),\ \delta x = {dx \over dt},\ \delta y = {dy \over dt},
    $$%
  to which we can add:
  $$%
    \delta x = {\partial x \over \partial t} (\delta t)
    \quad \text{and} \quad
    \delta t = 1 \ \Rightarrow \ \delta x = {dx \over dt}.
    $$%
  Of course,
  every vector $v$ can be realized as a variation of the point $x$ by considering $\gamma(t) = x +tv$.
\end{Article}

%%%%%%%%%%%%%%%%%%%%%%%%%%%%%%%%%%%%%%%%%%%%%%%%%%%%%%%%%%
%%
%% MARK: Higher Derivatives Components
%%
%%%%%%%%%%%%%%%%%%%%%%%%%%%%%%%%%%%%%%%%%%%%%%%%%%%%%%%%%%

\section{Higher derivatives components}

What has been done for the tangent linear map can be done for higher derivatives.
For example,
let $f \colon x \mapsto y$.
The second derivative
$$%
  \Der^2(f)(x) = \Der\big(x \mapsto \Der(f)(x)\big)(x)
  $$%
is represented itself by the bilinear form:
$$%
  \Der^2(f)(x)(v)(w) = \Der\big(x \mapsto \Der(f)(x)(v)\big)(x)(w),
  $$%
for all $v, w \in \RR^n$.
Let us decompose $v = \sum_{i=1}^n v^i \bme_i$ and $w = \sum_{j=1}^n w^j \bme_j$.
We get:
\begin{eqnarray*}
  \Der^2(f)(x)(v)(w) & = & \Der\bigg(x \mapsto \Der(f)(x)\bigg(\sum_{i=1}^n v^i \bme_i\bigg)\bigg)(x)\bigg(\sum_{j=1}^n w^j \bme_j\bigg) \\
  & = & \sum_{j=1}^n w^j \Der\big(x \mapsto \sum_{i=1}^n v^i \Der(f)(x)(\bme_i)\big)(x)(\bme_j) \\
  & = & \sum_{j=1}^n w^j \sum_{i=1}^n v^i  \Der\big(x \mapsto \Der(f)(x)(\bme_i)\big)(x)(\bme_j) \\
  & = & \sum_{j,i=1}^n w^j v^i {\partial \over \partial x^j}\left({\partial y \over \partial x^i}\right).
\end{eqnarray*}
And since $y = (y^1, \dots, y^m)$,
$$%
  \Der^2(f)(x)(v)(w) = \sum_{j,i=1}^n w^j v^i \Vector{
  {\partial^2 y^1 \over \partial x^j \partial x^i} \\[.1em] \vdots \\[.1em] {\partial^2 y^m \over \partial x^j \partial x^i}
  }
  = \sum_{j,i=1}^n w^j v^i \Vector{
  \partial^2_{ji} y^1 \\[.1em] \vdots \\[.1em] \partial^2_{ji} y^m
  }
  $$%
The partial derivatives $\partial_{ij} y^k$ are the components of the bilinear map $\Der^2(f)(x)$.

A main property of continuously differentiable functions is the commutativity of the partial derivatives:

\begin{Article}{Theorem (Schwarz).}{\itshape}
  Let $f \colon U \to \RR^m$,
  with $U$ a $n$-domain,
  be a $\Ck{2}$ map.
  Then,
  the second derivative $\Der^2(f)(x)$ is symmetric.
  In other words,
  the partial derivatives commute:
  $$%
    {\partial \over \partial x^i}\left({\partial y^k \over \partial x^j}\right) =
    {\partial \over \partial x^j}\left({\partial y^k \over \partial x^i}\right)
    \quad \text{i.e.} \quad \partial^2_{ij}y^k = \partial^2_{ji}y^k.
    $$%
\end{Article}
For $\Ck{k}$ maps,
the $k$-derivative $\Der^k(f)(x)$ is a $k$-linear map,
with components the partial derivatives:
$$%
  \partial^k_{i_1 \dots i_k} y^\ell  = \partial_{i_1} \big(\partial_{i_2 \dots i_k}^{k-1} y^\ell\big)
  $$%
And,
thanks to the last proposition,
this multilinear map is symmetric:
the order of the indices does not matter.

\begin{Article}{Theorem.}{\itshape}
  Let $f\colon U \to \RR^m$ be a map,
  where $U$ is an $n$-domain.
  The map $f$ is smooth if and only if $f$ admits partial derivatives at any order.
\end{Article}

\begin{proof}
  Indeed,
  and that is the main point:
  if all the partial derivatives at any order exist,
  that implies in particular that they are continuous,
  thanks to \textsection\,\ref{continuously-differentiable},
  and then the map $f$ is \emph{continuously infinitely differentiable}.
\end{proof}

This theorem is a practical criterion to check if a map is smooth.

That is everything we need to know to introduce and study diffeology.
That is what makes diffeology a good alternative to the usual teaching of differential geometry.
And we shall see in the future how it could be understood formally as a \emph{geometry in the sense of Felix Klein}.

%%%%%%%%%%%%%%%%%%%%%%%%%%%%%%%%%%%%%%%%%%%%%%%%%%%%%%%%%%
%%
%% MARK: The Category of Euclidean Domains
%%
%%%%%%%%%%%%%%%%%%%%%%%%%%%%%%%%%%%%%%%%%%%%%%%%%%%%%%%%%%

\section{The category of Euclidean domains}

Euclidean domains are the objects of a category for which the arrows are the smooth maps.
We denote it by \{Euclidean Domains\}.

The goal of diffeology is to transfer some properties of the category of Euclidean domains to arbitrary sets.

%%%%%%%%%%%%%%%%%%%%%%%%%%%%%%%%%%%%%%%%%%%%%%%%%%%%%%%%%%
%%
%% MARK: Some Theorems We Should Know
%%
%%%%%%%%%%%%%%%%%%%%%%%%%%%%%%%%%%%%%%%%%%%%%%%%%%%%%%%%%%

\section{Some theorems we should know}

There are a few important theorems of differential calculus in $\RR^n$ we will need in future development of diffeology,
for example the implicit function theorem or the rectification of vector fields and some others.
However,
they are not necessary for now to learn diffeology.
That is why it is better to introduce them only when they will be used.

\vfill
\centerline{\includegraphics{Figures/Notes-Beginning.pdf}}
